% Clase 16 — la integral de y dx sobre la curva proyectada C'. Se recorre por
% abajo (y_1) de x_min a x_max y se vuelve por arriba (y_2): al restar queda
% menos el área encerrada.
\begin{tikzpicture}[scale=1]

  % region
  \fill[figblue!12]
    plot[smooth cycle, tension=0.75] coordinates {(1,1.75) (2,0.8) (3.6,0.65) (4.9,1.35) (5,2.3) (4,3.05) (2.5,3.1) (1.3,2.6)};
  \draw[figline, line width=1pt]
    plot[smooth cycle, tension=0.75] coordinates {(1,1.75) (2,0.8) (3.6,0.65) (4.9,1.35) (5,2.3) (4,3.05) (2.5,3.1) (1.3,2.6)};

  % ejes
  \draw[figaxis] (-0.2,0) -- (6.2,0);
  \draw[figaxis] (0,-0.2) -- (0,3.9);
  \node[figlbl] at (6.45,0) {$x$};
  \node[figlbl] at (0,4.15) {$y$};

  % extremos
  \draw[figaux] (0.93,0) -- (0.93,2.0);
  \draw[figaux] (5.1,0) -- (5.1,1.9);
  \node[figlblaux] at (0.93,-0.32) {$x_{\min}$};
  \node[figlblaux] at (5.1,-0.32) {$x_{\max}$};

  % sentido de recorrido
  \draw[figvecres] (2.6,0.66) -- (3.3,0.62);
  \draw[figvecres] (3.4,3.13) -- (2.7,3.12);

  % rotulos
  \node[figlbl] at (3.0,0.28) {$y_1(x)$};
  \node[figlbl] at (3.05,3.5) {$y_2(x)$};
  \node[figlbl] at (3.0,1.9) {área de $C'$};
  \node[figlbl] at (5.55,2.75) {$C'$};

\end{tikzpicture}
